\chapter{Glossary of Doxonomic Terms}
\label{app:glossary}

\begin{description}[style=nextline, leftmargin=3cm]

\item[Anthropological regime] A dominant practical model of human nature specifying motivational theory, social ontology, and capacity claims. See Definition~\ref{def:anthropological-regime}.

\item[Asymptotic capture] The process by which sufficiently asymmetric evidence exposure drives beliefs to certainty regardless of prior. See Conjecture~\ref{conj:asymptotic-capture}.

\item[Belief production chain] The composition of operators from funding through institutional processing, narrative output, circulation, internalization, behavioral consequence, to reproduction. See Definition~\ref{def:bpc-formal}.

\item[Belief stock] The fraction of a population holding a given belief at a given time. See Definition~\ref{def:stock-flow}.

\item[Common sense] A proposition held with high credence and low variance across a population. See Definition~\ref{def:common-sense}.

\item[Common-sense capture] The phase transition by which a funded narrative becomes invisible as a claim, experienced as a feature of reality. See Definition~\ref{def:capture}.

\item[Conversion problem] The foundational problem of doxonomics: how resources are converted into durable public belief.

\item[Counter-doxonomics] The study and practice of resisting, disrupting, or replacing dominant doxonomic regimes. See Chapter~\ref{ch:counter-doxonomic}.

\item[Credibility collapse] Rapid decline in institutional credibility due to scandal, failure, or exposure. See Definition~\ref{def:credibility-collapse}.

\item[Doxa] (Greek) The unquestioned taken-for-granted; beliefs so deeply embedded they are not perceived as beliefs. See Definition~\ref{def:doxa}.

\item[Doxametry] The quantitative measurement branch of doxonomics. See Chapter~\ref{ch:doxometric-measurement}.

\item[Doxonomic equilibrium] A fixed point of the belief production chain where beliefs reproduce the conditions of their own production. See Definition~\ref{def:doxonomic-equilibrium}.

\item[Doxonomic externality] A social cost generated when a subsidized belief reshapes behavior in harmful ways. See Definition~\ref{def:doxonomic-externality}.

\item[Doxonomic game] An experimental economic game administered after exposure to a doxonomic treatment. See Definition~\ref{def:doxonomic-game}.

\item[Doxonomic network] A weighted directed graph of funders, institutions, media, and audiences connected by influence or resource flows. See Definition~\ref{def:doxonomic-network}.

\item[Doxonomic power] The capacity to shape the conditions under which people form beliefs, including the power to set priors, agendas, and epistemic rules. See Definition~\ref{def:doxonomic-power}.

\item[Doxonomic system] A structured ensemble of agents, institutions, narrative spaces, belief spaces, and operators for production and internalization. See Definition~\ref{def:doxonomic-system}.

\item[Doxonomics] The science of how material power shapes public belief. From Greek \textit{doxa} (opinion, belief) and \textit{-nomics} (systematic study).

\item[Epistemic broker] A node with high betweenness centrality in a doxonomic network, bridging otherwise disconnected narrative clusters. See Definition~\ref{def:broker}.

\item[Epistemic capital] Resources convertible to credibility, attention, or authority in belief formation. See Definition~\ref{def:epistemic-capital}.

\item[Epistemic democracy] A social arrangement in which belief production is not dominated by a narrow set of funders or institutions. See Definition~\ref{def:epistemic-democracy}.

\item[Epistemic fragmentation] A condition in which populations inhabit incommensurable belief worlds with no shared epistemic framework. See Definition~\ref{def:fragmentation}.

\item[Epistemic Herfindahl Index (EHI)] A concentration score measuring market share in belief production. See Definition~\ref{def:ehhi}.

\item[Epistemic rent] Excess doxonomic return earned from a monopoly position in legitimate knowledge production. See Definition~\ref{def:epistemic-rent}.

\item[Gluing obstruction] The failure of locally consistent belief systems to combine into a globally consistent whole, measured by sheaf cohomology. See Definition~\ref{def:obstruction}.

\item[Ideological ROI (IROI)] The belief change achieved per unit of doxonomic investment. See Definition~\ref{def:iroi}.

\item[Lock-in] A condition in which a belief is locally asymptotically stable with a large basin of attraction. See Definition~\ref{def:lock-in}.

\item[Narrative entropy] The information-theoretic measure of diversity in the narrative space. See Definition~\ref{def:narrative-entropy}.

\item[Narrative flow] The rate at which narrative content supporting a belief is produced and circulated. See Definition~\ref{def:stock-flow}.

\item[Narrative infrastructure] The durable network of institutions through which a society produces and reproduces its common sense. See Definition~\ref{def:narrative-infrastructure}.

\item[Narrative inoculation] Preemptive exposure to a weakened dominant narrative plus refutation, reducing subsequent susceptibility. See Definition~\ref{def:inoculation}.

\item[Narrative shield] A set of beliefs that deflects attention from structural causes to individual explanations. See Definition~\ref{def:narrative-shield}.

\item[Pluralistic ignorance] A condition in which most individuals privately doubt a proposition but publicly affirm it. See Definition~\ref{def:pluralistic-ignorance}.

\item[Prior-setting power] The capacity to influence the initial belief distribution of a population, typically through education. See Definition~\ref{def:prior-setting}.

\item[Tipping point] A discontinuous jump in equilibrium belief level caused by a small shift in threshold distribution. See Theorem~\ref{thm:tipping}.

\end{description}
